\chapter{Confronto sperimentale PIR vs mmWave}
\label{cap:confronto}

\begin{quote}\small
Questo capitolo risponde all'obiettivo O3. Riporta i risultati della campagna
sperimentale, svolta dal 18 agosto al 12 settembre 2026 su tre sensori, con 441 trial
validi e 31,9 ore di acquisizione. Protocollo, metriche e convenzioni sono quelli del
\cref{cap:metodologia}. Ogni risultato è accompagnato dalla sua dispersione su
ripetizioni indipendenti. Il rilevamento del respiro poggia sugli stessi dati ed è
trattato nel \cref{cap:vitalita} insieme all'indice di vitalità.
\end{quote}

\section{Stanza vuota: falsi positivi}
\label{sec:falsi-positivi}

L'acquisizione parte con l'operatore che esce subito dalla stanza. Si contano le
transizioni verso la presenza sul resto del file. Su 7,05~h di stanza vuota
(6,57~h di notte e 0,48~h di giorno) né il \ldb\ né il PIR hanno prodotto alcun
falso positivo. Con zero eventi osservati si riporta il limite superiore $3/T$ della
\cref{sec:metriche}. Il risultato è $\leq 0{,}43$ eventi/h per entrambi i sensori. La limitazione
di questo numero è il tempo di osservazione, non il sensore.

Il \ldc\ nella stessa stanza si comporta in modo molto diverso
(\cref{fig:falsi-positivi}). Con le soglie già tarate sul fondo misurato
(\cref{sec:ld2420-risultati}) produce 26,2 riaccensioni all'ora in una notte di 7~h.
La presenza è dichiarata nel 26\,\% del tempo senza nessuno nella stanza. Con la
successiva correzione del gate~0 scende a 5,4 eventi/h, dodici volte il limite
superiore del \ldb. Quest'ultimo valore viene però da un'ora di giorno e non da una
notte, quindi i due numeri non sono omogenei. Una seconda notte con la
configurazione corretta non è stata acquisita.

\begin{figure}[htbp]
\centering
\includegraphics[width=\linewidth]{figures/fig18_falsi_positivi}
\caption{Falsi positivi a stanza vuota, di notte, in scala logaritmica. \ldb\ a
soglie di fabbrica, zero eventi in 6,6~h, quindi il limite superiore $3/T$. \ldc\
con le soglie tarate sul fondo misurato, 26 riaccensioni/h in 7~h di notte. Con il
gate~0 portato sopra la coda del rumore, 6/h, misurati però in 1~h di giorno. I
conteggi della figura usano lo scarto di 60~s degli script e differiscono di poco da
quelli del testo, che usano la finestra del registro.}
\label{fig:falsi-positivi}

\end{figure}

\section{La persona ferma}
\label{sec:persona-ferma}

È lo scenario che risponde alla domanda di partenza. Il soggetto siede immobile a
2,30~m dai sensori, con la sola respirazione spontanea, per cinque trial da 302~s
utili a 27\,\textdegree C. I 7554 campioni con presenza reale danno il risultato della
\cref{tab:fermo-seduto}.

\begin{table}[htbp]
\centering
\caption{Persona ferma seduta a 2,30~m, cinque trial da 302~s.}
\label{tab:fermo-seduto}

\small
\begin{tabular}{@{}lrr@{}}
\toprule
 & \textbf{PIR \pir} & \textbf{Radar \ldb} \\
\midrule
Falsi negativi & $99{,}78 \pm 0{,}49\,\%$ & $\mathbf{0{,}00 \pm 0{,}00\,\%}$ \\
Trial con falsi negativi al 100\,\% & 4 su 5 & 0 su 5 \\
Campioni con presenza rilevata & $\sim$17 su 7554 & 7554 su 7554 \\
\bottomrule
\end{tabular}
\end{table}

In quattro trial su cinque il PIR ha prodotto zero rilevamenti in cinque minuti,
con una persona viva a 2,3~m in campo libero. Il radar ha rilevato la presenza nel
100\,\% dei campioni in tutti i trial. Non è un esito inatteso, perché è ciò che il
principio del piroelettrico prevede (\cref{sec:pir-principio}). La tesi lo
quantifica, con la sua dispersione, nella condizione che interessa il progetto. La
\cref{fig:timeline-sotto-banco} mostra l'andamento nel tempo di un trial in questa
geometria e di uno nella geometria del progetto (\cref{sec:sotto-banco}).

\begin{figure*}[htbp]
\forcerectofloat
\centering
\includegraphics[width=\linewidth]{figures/fig02_timeline_sotto_banco}
\caption{Presenza dichiarata dai due sensori su una persona immobile, in due
geometrie. A sinistra il soggetto seduto a 2,3~m, a destra rannicchiato sotto il
banco a circa 60~cm, cinque minuti ciascuno. Il radar è continuo, il PIR produce al
più un impulso isolato. La terza riga è la distanza riportata dal radar sul bersaglio
fermo, stabile entro 1,5~cm da seduto e molto più rumorosa sotto il piano. Gli
sbalzi a 200--300~cm sono il canale stazionario che per qualche campione si aggancia
a un riflesso lontano, con il canale in movimento e il PIR fermi. Non sono movimenti
del soggetto.}
\label{fig:timeline-sotto-banco}

\end{figure*}

\subsection{Che cosa misura il bit del PIR}
\label{sec:interpretazione-pir}
\label{sec:jumper}

La percentuale di campioni in cui il PIR è alto va letta sapendo come nasce. In
posizione L l'uscita è un monostabile a durata fissa. Su tutte le sessioni acquisite
in quella configurazione sono stati contati 195 impulsi con durata media
$3{,}46 \pm 0{,}09$~s, tutti fra 3,4 e 3,6~s e nessuno oltre 5~s. Vale anche nei trial
in cui il soggetto si muoveva per cinque minuti (\cref{fig:impulsi-pir}). La
percentuale è quindi il numero di eventi moltiplicato per la ritenuta, cioè un
conteggio di eventi travestito da frazione di tempo. In posizione H gli impulsi si
concatenano finché il movimento continua, fino a coprire un intero trial di 62~s. La
percentuale diventa allora una vera frazione di tempo con movimento rilevato. Il
significato delle due modalità è descritto nel \cref{sec:pir-principio}.

\begin{figure*}[htbp]
\forcerectofloat
\centering
\includegraphics[width=\linewidth]{figures/fig07_impulsi_pir}
\caption{Struttura degli impulsi del PIR. In configurazione L tutti gli impulsi
durano fra 3,4 e 3,8~s qualunque sia il movimento, perché l'uscita è un monostabile a
durata fissa. In configurazione H il ritrigger li concatena, fino a coprire l'intero
trial.}
\label{fig:impulsi-pir}

\end{figure*}

La prima parte della campagna è stata acquisita in posizione L. La posizione è stata
riconosciuta dai dati, perché una dispersione di 0,09~s su 195 impulsi è la firma di
una durata fissa. Il ponticello agisce solo dopo un trigger, quindi non può cambiare
gli scenari in cui il PIR non ha prodotto eventi. Le due serie con eventi ripetuti
(cammino sul posto a 1~m e micro-movimenti sotto il banco) sono state ripetute in H e
sono quelle riportate in questo capitolo. Per le altre l'insensibilità al ponticello è
stata verificata per misura, ripetendo in H il cammino a 2~m (0,00\,\% in entrambe le
posizioni) e l'immobilità sotto il banco (0,10\,\% in L, 1,32\,\% in H). Il dato
robusto, indipendente da ritenuta e ritrigger, è che sulla persona ferma il PIR non
ha prodotto eventi di movimento. Un evento che non esiste non può essere né allungato
né concatenato.

\section{Il tipo di movimento decide il PIR}
\label{sec:dose-risposta}
\label{sec:pir-movimento}
\label{sec:attraversamento}

Le due prove di questa sezione variano una sola grandezza alla volta. La prima varia
la quantità di movimento a distanza fissa, la seconda il tipo di movimento a parità
di distanza.

\paragraph{Dose-risposta a 1~m.} Soggetto in piedi a 1~m, ponticello in H, stesso
montaggio in tutte le condizioni. Cambia solo la quantità di movimento, su tre
livelli, con cinque trial per livello (\cref{tab:dose-risposta} e
\cref{fig:dose-risposta}). La distanza di 1~m è quella in cui il PIR funziona meglio,
quindi il confronto non è costruito a suo sfavore.

\begin{table*}[htbp]
\forceversofloat
\centering
\caption{Curva dose-risposta a 1~m, ponticello in H, cinque trial per condizione.
Fra le tre righe cambia soltanto la quantità di movimento del soggetto. Le due colonne
di destra sono grandezze del radar, riprese nel \cref{cap:vitalita}.}
\label{tab:dose-risposta}

\small
\begin{tabular}{@{}lrrrrr@{}}
\toprule
\textbf{Condizione} & \textbf{PIR} & \textbf{dev.\ rel.} & \textbf{Radar} &
\textbf{Energia mov.} & \textbf{Disp.\ dist.} \\
\midrule
Immobile          & $1{,}52 \pm 1{,}03\,\%$   & $68\,\%$ & $\mathbf{100\,\%}$ & 78,1 & 20,2~cm \\
Micro-movimenti   & $51{,}60 \pm 21{,}72\,\%$ & $\mathbf{42\,\%}$ & $\mathbf{100\,\%}$ & 86,1 & 18,3~cm \\
Cammino sul posto & $85{,}20 \pm 11{,}52\,\%$ & $14\,\%$ & $\mathbf{100\,\%}$ & 98,2 & 11,0~cm \\
\bottomrule
\end{tabular}
\end{table*}

\begin{figure*}[htbp]
\forceversofloat
\centering
\includegraphics[width=\linewidth]{figures/fig01_dose_risposta}
\caption{La curva dose-risposta a 1~m. A geometria, distanza e configurazione
costanti cambia solo la quantità di movimento. Il PIR passa dall'1,5 all'85\,\%,
con la dispersione massima nella zona intermedia. Il radar resta al 100\,\% in tutte
e tre le condizioni e la sua energia in movimento cresce in modo monotono.}
\label{fig:dose-risposta}

\end{figure*}

Il PIR è monotono nella quantità di movimento, dall'1,52 all'85,20\,\%. Il radar
resta al 100\,\% in tutte le condizioni. Fra le righe cambia una sola variabile,
quindi il fallimento del PIR sulla persona ferma non dipende dalla distanza (a 1~m
funziona bene), né dalla configurazione (è in H), né dalla taratura (invariata).
Dipende dall'immobilità del soggetto. La struttura degli impulsi lo rende concreto.
In cammino il PIR produce impulsi lunghi e concatenati, sei impulsi completi da 14,1~s
in media e cinque troncati dal bordo della finestra, il più lungo dei quali copre
l'intero trial. Da immobile produce due impulsi in 1010~s di acquisizione.

Il dato più rilevante per un'applicazione di soccorso è la dispersione della
condizione intermedia, $\pm 21{,}72$ punti su una media di 51,60. Due trial identici
per l'operatore danno risultati molto diversi. Vicino alla propria soglia il PIR non è
soltanto meno sensibile, è imprevedibile. Una prestazione imprevedibile è peggiore di
un limite noto, perché non permette di interpretare l'assenza di segnale.

\paragraph{Sul posto contro attraversamento.} Il cammino sul posto è stato ripetuto
a 1, 2, 3, 4 e 5~m, con cinque trial per distanza. A 1~m il PIR rileva l'85,2\,\% del
tempo. A 2~m il tasso è esattamente zero e resta zero a 3, 4 e 5~m, venti trial
consecutivi al 100\,\% di falsi negativi con un soggetto in movimento continuo
davanti al sensore. Il crollo fra 1 e 2~m non è graduale. Questo esito è in aperto
contrasto con la portata di 3--7~m dichiarata dal costruttore. Da solo non distingue
fra due spiegazioni, il tipo di movimento oppure un esemplare guasto o tarato male.
Per decidere, il soggetto ha ripetuto la prova camminando in modo continuo
trasversalmente all'asse del sensore, a distanza costante. Tutto il resto è
invariato, ponticello in H e trimmer di sensibilità nella stessa posizione, con tre
trial per distanza. Il risultato è in \cref{fig:tipo-movimento}.

\begin{figure*}[htbp]
\forcerectofloat
\centering
\includegraphics[width=\linewidth]{figures/fig28_tipo_movimento}
\caption{Il tipo di movimento decide il PIR. A parità di distanza, configurazione e
sensibilità, il cammino sul posto è rilevato solo a 1~m (85\,\%) e mai da 2~m in su.
L'attraversamento trasversale è rilevato al 99--100\,\% da 2 a 5~m. Il radar è al
100\,\% con entrambi i movimenti. A 3--5~m le serie sul posto sono in posizione L,
ma contengono zero eventi e non dipendono dal ponticello.}
\label{fig:tipo-movimento}

\end{figure*}

Un sensore che rileva un attraversamento a 5~m nel 98,7\,\% dei campioni non è né
guasto né poco sensibile. La portata dichiarata risulta confermata. La variabile che
decide se il PIR rileva non è la distanza, è il tipo di movimento. A qualunque
distanza fra 2 e 5~m il tasso passa da zero a circa il 99\,\% cambiando soltanto la
direzione dello spostamento. È il comportamento che la geometria della lente di
Fresnel prevede (\cref{sec:fresnel}), qui misurato. Una prova aggiuntiva a 1~m nella
stessa sessione lo conferma a stimolo radar costante. L'energia in movimento del
radar sta fra 96,5 e 98,1 in tutte le prove. Il PIR rileva il 2\,\% con il movimento
delle braccia e il 100\,\% con l'oscillazione laterale del busto da seduto, che è una
traversata in miniatura. I due sensori non misurano la stessa grandezza con
sensibilità diverse. Misurano due grandezze diverse.

Per il progetto la conseguenza è diretta. Il movimento per cui il PIR è progettato,
attraversare un varco o percorrere un corridoio, è l'opposto di quello di una persona
intrappolata sotto un arredo, che si muove sul posto quando si muove. Il limite non è
di portata, è strutturale rispetto allo scenario.

\paragraph{Il radar fornisce indicatori graduati.} Le ultime due colonne della
\cref{tab:dose-risposta} sono un risultato a sé. L'energia del bersaglio in movimento
cresce con la quantità di movimento (78,1, 86,1, 98,2). La dispersione della distanza
cala (20,2, 18,3, 11,0~cm), perché un'eco più forte dà una stima più stabile. Sono due
grandezze continue e monotone misurate a geometria costante. Su di esse poggia
l'indice di vitalità del \cref{cap:vitalita}.

\section{Lo scenario DIPME: la persona sotto il banco}
\label{sec:sotto-banco}
\label{sec:gate-sotto-banco}

È lo scenario che replica la condizione del progetto. I sensori sono fissati sotto il
piano di un tavolo e il soggetto è rannicchiato sotto di esso a 50--60~cm
(\cref{fig:geometrie}b). Cinque trial da 302~s utili in due condizioni, immobilità e
presenza di micro-movimenti, entrambe con il ponticello del PIR in H. Il \ldc\ ha
ripetuto le stesse condizioni nella stessa geometria, e con esse anche le due
condizioni a 1~m in piedi. La riga a tre sensori è così completa
(\cref{fig:uprise-tre-sensori}).

\begin{figure*}[htbp]
\forceversofloat
\centering
\includegraphics[width=\linewidth]{figures/fig19_uprise_tre_sensori}
\caption{Lo scenario DIPME a tre sensori, cinque trial per condizione. PIR e \ldb\
in agosto, \ldc\ in settembre nella stessa geometria, con soglie tarate e scarto di
90~s. I due radar sono al 100\,\% in tutte le condizioni, salvo il \ldc\ al 98,4\,\%
da immobile sotto il banco per un rilascio di 24~s in un trial su cinque. Il PIR
passa dall'1--2\,\% da fermo al 52--97\,\% con i micro-movimenti.}
\label{fig:uprise-tre-sensori}

\end{figure*}

Sotto il banco il contrasto è ancora più netto che a 1~m. Con i micro-movimenti il PIR
rileva il $96{,}52 \pm 3{,}60\,\%$ del tempo, da immobile l'$1{,}32 \pm 0{,}86\,\%$.
A 60~cm il PIR è un rilevatore di movimento quasi perfetto e resta cieco alla persona
ferma. Il \ldb\ è al 100\,\% in entrambe le condizioni su 8054 campioni per
condizione. Il \ldc\ è al 100\,\% con i micro-movimenti e al $98{,}4 \pm 3{,}6\,\%$
da immobile, con un solo rilascio di 24~s in un trial. È il primo falso negativo su
persona presente della campagna del secondo radar, contro zero del \ldb\ in venti
trial. Insieme alla dose-risposta di \cref{sec:dose-risposta} questa è la seconda
coppia appaiata del capitolo. Due distanze, due geometrie, la stessa conclusione.

\paragraph{Il radar rileva bene ma localizza male.} La dispersione della distanza
entro il trial vale 14--24~cm sotto il banco, contro $1{,}48 \pm 0{,}77$~cm sul
soggetto seduto a 2,3~m (\cref{fig:timeline-sotto-banco}). È circa dieci volte peggio
su un bersaglio immobile più vicino. Per il progetto è irrilevante, perché al
soccorritore serve sapere se c'è qualcuno sotto quel banco e non a quanti centimetri.
Esclude però l'uso della distanza come misura fine in questa geometria. Gli sbalzi
verso 200--300~cm visibili in \cref{fig:timeline-sotto-banco} riguardano dallo 0,3 al
7,4\,\% dei campioni nei cinque trial. Avvengono con il canale in movimento meno
attivo del solito e con il PIR a zero, quindi non sono movimenti del soggetto. È il
canale stazionario che si aggancia a un riflesso lontano. I file terminano per durata
prefissata prima che l'operatore si alzi, e i primi 40~s di posizionamento sono
scartati. La lettura per gate spiega il perché. Sul \ldb\ i gate stazionari 0 e 1
sono strutturalmente nulli (\cref{sec:senergy-zero}) e la presenza è retta dai gate 2
e 3, con un'energia che resta a 24 anche al gate~8 contro 6,5 nello scenario da
seduto. Sul \ldc\ la presenza da immobile è retta dal gate~2 (70--140~cm) e non dal
gate~1 in cui la persona sta. Lo stesso fenomeno compare su due moduli diversi,
verosimilmente per cammini multipli sotto il piano. È una spiegazione coerente, non
una misura, perché il segnale grezzo non è accessibile. Il punto operativo è comunque
acquisito. Sotto il metro l'indice di vitalità deve usare i canali in movimento
(\cref{cap:vitalita}).

\section{Accuratezza della distanza ed energia}
\label{sec:distanza}
\label{sec:energia-distanza}

I 25 trial di cammino sul posto a 1--5~m caratterizzano la misura di distanza del
\ldb\ (\cref{tab:distanza}, \cref{fig:distanza-regressione}).

\begin{table*}[htbp]
\forcerectofloat
\centering
\caption{Accuratezza della distanza riportata dal \ldb, 25 trial (5 distanze $\times$
5 ripetizioni, 62~s utili ciascuno). ``Misurata'' è la media fra trial con la sua
deviazione standard, ``residuo'' lo scarto dalla retta dell'\cref{eq:regressione},
``disp.\ entro trial'' la deviazione standard della distanza dentro un trial, mediata
sui cinque. L'energia è il valore medio del bersaglio in movimento.}
\label{tab:distanza}

\small
\begin{tabular}{@{}rrrrrrr@{}}
\toprule
\textbf{Reale} & \textbf{Misurata} & \textbf{Errore} & \textbf{Err.\ rel.} &
\textbf{Residuo} & \textbf{Disp.\ entro trial} & \textbf{Energia} \\
\midrule
1~m & $99{,}48 \pm 2{,}61$~cm  & $-0{,}52$~cm & $-0{,}5\,\%$ & $-3{,}0$~cm & $10{,}4$~cm & 99,0 \\
2~m & $211{,}22 \pm 0{,}77$~cm & $+11{,}22$~cm & $+5{,}6\,\%$ & $+4{,}9$~cm & $14{,}3$~cm & 85,1 \\
3~m & $309{,}76 \pm 1{,}13$~cm & $+9{,}76$~cm & $+3{,}3\,\%$ & $-0{,}4$~cm & $17{,}5$~cm & 54,4 \\
4~m & $411{,}84 \pm 2{,}76$~cm & $+11{,}84$~cm & $+3{,}0\,\%$ & $-2{,}1$~cm & $26{,}4$~cm & 35,1 \\
5~m & $518{,}20 \pm 2{,}13$~cm & $+18{,}20$~cm & $+3{,}6\,\%$ & $+0{,}5$~cm & $11{,}7$~cm & 27,6 \\
\bottomrule
\end{tabular}
\end{table*}

\begin{figure*}[htbp]
\forceversofloat
\centering
\includegraphics[width=\linewidth]{figures/fig04_distanza_regressione}
\caption{Accuratezza della distanza del \ldb, 25 trial fra 1 e 5~m. La retta è
calcolata sulle cinque medie per distanza ($R^2 = 0{,}99965$). Sui 25 trial singoli si
otterrebbe 0,99949. L'errore di scala del $+3{,}81\,\%$ ha offset praticamente nullo
ed è correggibile con un solo coefficiente.}
\label{fig:distanza-regressione}

\end{figure*}

La regressione lineare fra distanza misurata e reale dà
\begin{equation}
d_{\text{misurata}} = 1{,}0381 \cdot d_{\text{reale}} - 1{,}32~\text{cm},
\qquad R^2 = 0{,}99965.
\label{eq:regressione}
\end{equation}
Il radar è estremamente lineare, con residui dalla retta non superiori a 4,9~cm su
tutto il range. L'errore è quasi interamente di scala ($+3{,}81\,\%$) con offset
trascurabile. Questo spiega perché l'errore grezzo cresce con la distanza, da
$-0{,}5$~cm a 1~m a $+18{,}2$~cm a 5~m. È correggibile dividendo la lettura per
1,0381, dopo di che il residuo scende sotto i 5~cm a tutte le distanze. I dati non
permettono di attribuire il $+3{,}81\,\%$ al modulo piuttosto che al riferimento con
cui sono state segnate le distanze sul pavimento. Servirebbe un distanziometro
indipendente.

La dispersione entro il trial (10--26~cm) non è rumore del sensore. È l'oscillazione
reale del busto di chi cammina sul posto, come mostra il confronto con l'1,48~cm del
soggetto seduto immobile. In quello scenario la variabilità fra trial della distanza
media ($\pm 2{,}38$~cm) è maggiore di quella entro il trial. Il fattore limitante è
come il soggetto si siede, non lo strumento.

\begin{figure}[htbp]
\centering
\includegraphics[width=\linewidth]{figures/fig05_energia_distanza}
\caption{Energia del bersaglio in movimento e dispersione della distanza entro il
trial in funzione della distanza, 25 trial. Il decadimento dell'energia è molto più
lento di quello previsto dall'equazione del radar, perché il valore è un indice
normalizzato 0--100 e non una potenza. A 5~m l'energia (27,6) è ancora quasi il
doppio della soglia di fabbrica dei gate lontani.}
\label{fig:energia-portata}

\end{figure}

L'energia media del bersaglio decade da 99 a 1~m a 28 a 5~m
(\cref{fig:energia-portata}). Il decadimento è molto più lento della dipendenza
$1/D^4$ della potenza ricevuta da un bersaglio puntiforme. Il valore è un indice
normalizzato prodotto da un'elaborazione interna non documentata. Se ne usa la
monotonia e non il valore assoluto in termini fisici. Il range di 1--5~m è un limite
della stanza, non del sensore. La portata oltre i 5~m è misurata in corridoio in
\cref{sec:portata-massima}.

\paragraph{Il secondo radar e il PIR.} Il PIR non compare in questa sezione perché
non misura né una distanza né un'energia. Il suo dato è un bit. Il \ldc\ riporta una
distanza solo sul bersaglio in movimento e, sull'esemplare in dotazione, solo da
vicino. In venti trial di cammino sul posto da 0,5 a 2~m la distanza riportata vale
$59 \pm 5$~cm a 0,5~m e $121 \pm 9$~cm a 1~m, con errori di $+9$ e $+21$~cm e da 4 a
38 valori distinti al minuto. A 1,5 e 2~m il modulo riporta un solo valore stantio per
trial, quindi una retta di calibrazione non è costruibile. Sulla persona immobile
riporta il valore di riposo (2--5~cm) con la presenza al 100\,\%, come dichiara il
produttore. Anche l'energia del gate della persona non è una grandezza graduata su
questo esemplare, come mostra la dose-risposta ripetuta in \cref{sec:ld2420-risultati}.

\section{Latenze di rilevamento e di rilascio}
\label{sec:latenza}
\label{sec:rilascio}

L'acquisizione parte con la stanza vuota, oppure con il soggetto presente. Il
soggetto entra, oppure esce, a $t = 30$~s, istante annunciato a voce dallo script di
acquisizione. Lo script conta dalla prima riga di dati ricevuta, cioè dalla stessa
origine dei tempi dello script di analisi. Un cronometro esterno non sarebbe adeguato,
perché fra l'apertura della porta seriale e la prima riga passano 2--3~s. Dieci trial
all'ingresso e cinque all'uscita, tutti validi, con nessun sensore attivo prima
dell'evento. I risultati sono in \cref{tab:latenze} e \cref{fig:latenze}.

\begin{table*}[htbp]
\forceversofloat
\centering
\caption{Latenze all'ingresso e tempo di rilascio all'uscita. \ldb\ e PIR letti
nello stesso trial (10 all'ingresso, 5 all'uscita), con la differenza appaiata
calcolata trial per trial e nessuna riaccensione dopo il rilascio. Il \ldc\ è misurato
in una sessione separata (5 trial all'ingresso, 3 all'uscita) con il ritardo di
scomparsa portato da 30 a 5~s come sul \ldb. Il PIR non era cablato, quindi manca la
differenza appaiata.}
\label{tab:latenze}

\small
\begin{tabularx}{\linewidth}{@{}lrrrX@{}}
\toprule
\textbf{Grandezza} & \textbf{\ldb} & \textbf{PIR \pir} & \textbf{\ldc} & \textbf{Nota} \\
\midrule
Latenza all'ingresso & $5{,}36 \pm 0{,}30$~s & $5{,}96 \pm 0{,}31$~s & $7{,}28 \pm 1{,}40$~s &
  latenza operativa, non del solo sensore \\
Differenza appaiata \ldb\ $-$ PIR & \multicolumn{2}{c}{$\mathbf{-0{,}60 \pm 0{,}35}$~\textbf{s}} & --- &
  il radar rileva per primo in 10 trial su 10 \\
Tempo di rilascio all'uscita & $\mathbf{18{,}36 \pm 0{,}54}$~s & $12{,}96 \pm 1{,}40$~s & $7{,}7 \pm 0{,}5$~s &
  \ldb\ più lento del PIR di $5{,}40 \pm 0{,}67$~s. \ldc\ con ritardo 5~s \\
\bottomrule
\end{tabularx}
\end{table*}

\begin{figure*}[htbp]
\forceversofloat
\centering
\includegraphics[width=\linewidth]{figures/fig06_latenze}
\caption{A sinistra le latenze appaiate all'ingresso, con il radar primo in dieci
trial su dieci. A destra la latenza di rilascio all'uscita, dove il radar impiega 18~s
a dichiarare la stanza vuota contro un timeout configurato di 5.}
\label{fig:latenze}

\end{figure*}

\paragraph{Ingresso.} Le latenze assolute intorno a 5--6~s sono dominate dal tragitto
di rientro del soggetto e dal suo tempo di reazione all'annuncio. Vanno lette come
latenza operativa, limite superiore di quella del sensore. La grandezza pulita è la
differenza appaiata, perché tragitto e reazione sono identici per i due sensori nello
stesso trial e si cancellano. Il valore $-0{,}60 \pm 0{,}35$~s è una proprietà dei
sensori, con errore standard della media di 0,11~s. Il segno è concorde in tutti e
dieci i trial, con probabilità $\approx 0{,}002$ sotto l'ipotesi di equivalenza. Con
il campionamento a 5~Hz le singole differenze valgono da uno a sette campioni. La
conclusione poggia sulla concordanza del segno, non sull'ampiezza di una misura. Ci si
attendeva il contrario, perché attraversare una porta è lo stimolo trasversale per
cui la lente di Fresnel è progettata. Una spiegazione plausibile, non dimostrata, è
che il radar agganci il soggetto ancora in avvicinamento grazie ai 6~m di portata.

\paragraph{Uscita.} Qui il radar è nettamente più lento, $5{,}40 \pm 0{,}67$~s in
più del PIR, un effetto otto volte la propria incertezza già con cinque trial. Il
tempo di rilascio include il tempo di uscita dal campo, e il PIR può servire da
cronometro. La sua uscita ricade $3{,}46$~s dopo l'ultimo trigger. Il soggetto esce
quindi dal campo circa 9,5~s dopo l'annuncio, e la coda propria del radar vale circa
$18{,}36 - 9{,}5 \approx 8{,}9$~s. È una stima, perché i campi dei due sensori non
coincidono, ma concorda con altre due osservazioni indipendenti (10 e 12~s). Il
timeout di presenza letto dal modulo vale 5~s (\cref{tab:identita}). Il parametro
configurato non descrive quindi il comportamento osservato. Per il progetto una
stanza appena liberata risulta occupata per quasi dieci secondi.

\paragraph{Il secondo radar.} Sullo stesso percorso il \ldc\ rileva l'ingresso in
$7{,}28 \pm 1{,}40$~s su cinque trial, quasi due secondi dopo il \ldb\ e con una
dispersione quattro volte maggiore. La ragione è che aggancia il soggetto solo entro
circa 1~m, mentre il \ldb\ lo aggancia in avvicinamento a 4--5~m. Il rilascio va
letto insieme al suo parametro. Con il ritardo di scomparsa di fabbrica, 30~s, il
modulo restava in presenza per 87--120~s dopo l'uscita, perché basta un superamento
della soglia di mantenimento ogni mezzo minuto a riarmare il ritardo
(\cref{sec:ld2420-risultati}). Portato a 5~s come sul \ldb, il rilascio scende a
$7{,}7 \pm 0{,}5$~s. Il confronto a parametro pari (7,7 contro 18,4~s) è dominato dal
tempo di uscita dal campo, 1--2~m di portata contro 5--6. Non significa quindi che il
\ldc\ rilasci più in fretta.

\section{Copertura angolare e selettività fra banchi}
\label{sec:angolare}
\label{sec:selettivita}

Nell'architettura del progetto ogni arredo ha il proprio sensore, che deve riportare
il proprio occupante e non quello del banco accanto. Lo strumento offerto dal \ldb\ è
il gate massimo, che con la regola documentata di 75~cm per gate porta la portata a
150~cm quando vale 2 (\cref{sec:soglie}). Il gate limita però la distanza e non
l'angolo, quindi la domanda è doppia. Il limite di distanza funziona? Il vicino
laterale è visto?

\paragraph{Selettività con gate massimo 2.} Quattro scenari con soggetti sempre
fermi, tre trial da 182~s ciascuno (\cref{fig:selettivita}). L'occupante a 1~m
sull'asse è rilevato al 100\,\%. Una persona a 3~m sull'asse, oltre il taglio, è
rilevata allo 0,00\,\%. Una persona a 1~m a 90\textdegree\ è rilevata allo 0,00\,\% a
regime. Il valore $24{,}8 \pm 19{,}6\,\%$ che si ottiene con lo scarto standard di
20~s è interamente la coda del posizionamento, un tratto continuo dall'istante zero
che si spegne una sola volta (dopo 66, 29 e 101~s nei tre trial) e non si riaccende
più. Con l'occupante e il vicino insieme, nessuna grandezza dell'occupante cambia
oltre $2\,\sigma$ (la distanza di 4~cm, l'energia di 4 punti, con tre trial per
scenario).

\begin{figure*}[htbp]
\forceversofloat
\centering
\includegraphics[width=\linewidth]{figures/fig09_selettivita}
\caption{Selettività spaziale con gate massimo 2 (150~cm), tre trial per scenario.
Le barre chiare usano lo scarto ordinario dei primi 20~s, quelle scure lo scarto di
120~s (a regime). A 3~m zero rilevamenti. Il vicino a 90\textdegree\ fermo non è
rilevato a regime, e il 24,8\,\% della finestra ordinaria è la coda del
posizionamento, che in questi scenari dura fino a 101~s.}
\label{fig:selettivita}

\end{figure*}

\paragraph{Copertura angolare.} Il soggetto siede su un arco di raggio 1~m a sei
azimut (0, 45, 60, 75, 90 e 120\textdegree), con la sedia rivolta verso i sensori a
ogni posizione, e oscilla lateralmente il busto con ampiezza costante. Gate massimo 2,
da uno a tre trial per angolo (\cref{fig:angolare-due-radar}). Il \ldb\ rileva al
100\,\% fino a 90\textdegree\ compresi, con l'energia che cala di due punti soli fra
l'asse e i 90\textdegree\ (da 96,5 a 94,3), e crolla fra 90 e 120\textdegree. La
copertura eccede nettamente i $\pm 60$\textdegree\ del manuale. Il PIR collassa prima,
dal 100\,\% sull'asse al 12\,\% a 90\textdegree\ e a zero a 120\textdegree. Fra 45 e
75\textdegree\ la dispersione fra trial è però dell'ordine dell'effetto, e a
45\textdegree\ i tre trial danno 100, 100 e 33\,\%. Le sole affermazioni sostenibili
sono i tre estremi. L'irripetibilità è essa stessa un risultato, la terza occorrenza
dopo la zona intermedia della dose-risposta. Il \ldc\ a 1~m ha il fascio utile fino a
circa 75\textdegree\ ed è al fondo a 90\textdegree, più stretto del \ldb\ e più largo
dei $\pm 45$ o $\pm 60$\textdegree\ del suo manuale. Il confine è coerente anche con
il margine ridotto del trasmettitore dell'esemplare (\cref{sec:ld2420-stato}).

\begin{figure*}[htbp]
\forcerectofloat
\centering
\includegraphics[width=\linewidth]{figures/fig20_angolare}
\vspace{3.0cm}% spazio per la didascalia nel margine (tufte non lo riserva)
\caption{Copertura angolare a 1~m, busto laterale da seduto. A: il \ldb\ rileva al
100\,\% fino a 90\textdegree\ e crolla a 120, ben oltre i $\pm 60$\textdegree\
dichiarati. Il PIR non è ripetibile fra 45 e 75\textdegree. B: l'esemplare di \ldc\
è al fondo a 90\textdegree, fascio utile fino a circa 75.}
\label{fig:angolare-due-radar}

\end{figure*}

\paragraph{La selettività per configurazione non è ottenibile.} Incrociando le due
prove alla stessa distanza, allo stesso angolo e con lo stesso gate, il vicino a 1~m a
90\textdegree\ è rilevato allo 0\,\% se è fermo e al 100\,\% se si muove. Ciò che
nella prima prova sembrava selettività era la debolezza dell'eco di un corpo immobile.
Ridurre il gate non protegge, perché limita la distanza e non l'angolo. Il diagramma
di irradiazione non attenua abbastanza a 90\textdegree\ da escludere un bersaglio in
movimento. Un banco vuoto accanto a una persona che si muove risulta occupato. La
separazione fra arredi adiacenti va cercata nel montaggio fisico (orientamento verso
il basso, schermatura del lobo laterale) e non nei parametri del modulo. Fuori asse il
\ldb\ sottostima anche la distanza, 108, 107, 95, 94 e 85~cm da 0 a 90\textdegree\
sull'arco a 100~cm, verosimilmente agganciando la porzione di corpo più vicina.

\section{Due persone davanti al radar}
\label{sec:due-persone}

Il \ldb\ non ha risoluzione angolare, ma espone due canali indipendenti, uno per il
bersaglio in movimento e uno per quello stazionario, ciascuno con la propria distanza
(\cref{sec:protocollo-seriale}). Persona A ferma a 2~m e persona B a 4~m, tre trial da
102~s per geometria, sfalsate lateralmente di 50~cm per evitare l'ombra reciproca. La
variante in fila è acquisita come confronto. Un trial di controllo con B da sola
ferma a 4~m dà rilevamento al 100\,\% e distanza $388{,}9 \pm 3{,}7$~cm, quindi il
confondente della distanza è escluso (\cref{tab:due-persone}, \cref{fig:due-persone}).

\begin{table*}[htbp]
\forcerectofloat
\centering
\caption{Capacità di separare due persone, su tre geometrie. ``B riportata'' è la
frazione di campioni in cui la persona lontana compare su almeno uno dei due canali,
riconosciuta dalla banda di distanza. Tre trial per geometria, 1833 campioni utili
ciascuna.}
\label{tab:due-persone}

\small
\begin{tabular}{@{}lrrrrr@{}}
\toprule
\textbf{Geometria} & \textbf{Entrambi i} & \textbf{B} & \textbf{Dist.} &
\textbf{Dist.} & \textbf{Presenza} \\
 & \textbf{canali attivi} & \textbf{riportata} & \textbf{staz.} & \textbf{mov.} &
\textbf{radar} \\
\midrule
B cammina, sfalsate di lato & $79{,}2\,\%$ & $\mathbf{73{,}8\,\%}$ & 212~cm & 389~cm & $100\,\%$ \\
Entrambe ferme, sfalsate    & $21{,}2\,\%$ & $\mathbf{19{,}5\,\%}$ & 220~cm & 375~cm & $100\,\%$ \\
B cammina, in fila dietro A & $42{,}4\,\%$ & $\mathbf{0{,}3\,\%}$  & 215~cm & 215~cm & $100\,\%$ \\
\bottomrule
\end{tabular}
\end{table*}

\begin{figure*}[htbp]
\forceversofloat
\centering
\includegraphics[width=\linewidth]{figures/fig08_due_persone}
\caption{Due persone davanti al radar in tre geometrie. Il modulo riporta un
bersaglio per canale, quindi separa bene due persone in stati diversi, male due
persone ferme, per nulla una persona dietro l'altra. La presenza resta al 100\,\% in
ogni caso.}
\label{fig:due-persone}

\end{figure*}

Il limite del modulo va formulato così, riporta un bersaglio per canale. Con A ferma
e B in movimento i due canali sono attivi insieme nel 79\,\% dei campioni e riportano
due distanze separate da 177~cm. L'energia in movimento (38,4) è coerente con quella
misurata a 4~m nella caratterizzazione della distanza (35,1). Con entrambe ferme i
due bersagli competono per lo stesso canale e la più vicina vince nel 93\,\% dei
campioni. B compare comunque nel 19,5\,\% grazie ai suoi micro-movimenti involontari.
In fila il corpo davanti nasconde del tutto quello dietro, e i due canali riportano lo
stesso bersaglio, con differenza mediana di 3~cm. In ogni geometria la presenza è al
100\,\%. Il radar non sbaglia mai a dire che c'è qualcuno, perde il conteggio. Per il
progetto questo rafforza l'architettura a un sensore per arredo, dove ogni sensore ha
un solo bersaglio. Esclude invece l'uso di un sensore singolo come contatore di
presenze in aula.

\paragraph{Il secondo radar.} Il \ldc\ ha un solo canale di bersaglio
(\cref{sec:ld2420}), e la prova ripetuta con A ferma a 1~m e B a 2~m lo mostra. La
presenza è al 100\,\% in tutte le geometrie, ma non esiste un canale su cui B possa
comparire. In fila B è invisibile come sul \ldb. Sfalsate, le energie per gate non
separano le due persone. La distanza riportata si aggancia a B quando B è in
movimento (211--214~cm in due trial su tre) e resta piantata su un valore fisso
(268~cm) con entrambe ferme. Il secondo radar dice che c'è qualcuno e non quanti,
anche nella geometria in cui il \ldb\ riporta la seconda persona nel 74\,\% dei
campioni.

\section{Penetrazione degli ostacoli}
\label{sec:ostacoli}
\label{sec:ostacoli-portata}
\label{sec:ld2420-attenuazione}

Il pannello è interposto a 20~cm dal sensore, perpendicolare all'asse e libero di
reggersi da solo. Serve un soggetto in movimento sul posto, perché su un bersaglio
fermo il canale stazionario satura e l'attenuazione non sarebbe osservabile. Servono
anche due distanze, 3~m per il \ldb, dove l'energia vale circa metà scala, e 1~m per
il PIR, che a 3~m non rileva il movimento sul posto. Il \ldc\ ha ripetuto la serie a
1~m, l'unica distanza a cui l'esemplare lavora. Le sue energie sono interi a 16~bit
che lo strumento del produttore mostra in decibel, quindi la sua attenuazione è
espressa come $10\log_{10}(E_0/E)$ sul gate della persona. È un'ipotesi dichiarata,
perché il produttore non documenta il campo. Il riferimento senza ostacolo è la media
di sei trial a inizio e fine sessione, perché la baseline è la misura meno
riproducibile della serie (da 47,8 a 57,3 sul \ldb). I risultati sono in
\cref{tab:ostacoli} e \cref{fig:attenuazione}.

\begin{table*}[htbp]
\forcerectofloat
\centering
\caption{Attenuazione per materiale sui due radar e rilevamento del PIR a 1~m. \ldb:
riduzione dell'energia del bersaglio in movimento a 3~m rispetto alla baseline
mediata (53,45), in unità dello strumento. \ldc: attenuazione in decibel del gate~2
grezzo a 1~m (fondo 13). Legno e metallo sono al fondo e quindi limiti inferiori,
perché la dinamica dell'esemplare è di 5~dB. I pannelli hanno coperture angolari
diverse ($\pm 42$ a $\pm 70$\textdegree), quindi le attenuazioni sono limiti
inferiori. Il confronto più pulito è legno contro cartone, entrambi a
$\pm 54$\textdegree.}
\label{tab:ostacoli}

\small
\begin{tabular}{@{}llrrr@{}}
\toprule
\textbf{Materiale} & \textbf{Spessore} & \textbf{\ldb\ (3~m)} & \textbf{\ldc\ (1~m)} & \textbf{PIR (1~m)} \\
\midrule
Nessuno (riferimento) & --- & --- & --- & $79{,}6\,\%$ \\
Plastica (tanica cava) & $2 \times 5$~mm + aria & $-7{,}5\,\%$ & 0,4~dB & $0\,\%$ \\
Cartongesso & 10~mm & $-26\,\%$ & \textbf{0,7~dB} & $0\,\%$ \\
Cartone (scatola) & $2 \times 5$~mm + aria & $-17{,}2\,\%$ & 1,0~dB & $0\,\%$ \\
Vetroresina & 1~mm & $-19{,}6\,\%$ & 1,0~dB & $0\,\%$ \\
Vetro & 5~mm & $-40{,}9\,\%$ & 2,4~dB & $\sim 0\,\%$ \\
Legno & 10~mm & $-41{,}6\,\%$ & $\geq 4{,}1$~dB & $\sim 0\,\%$ \\
Metallo & 1~mm & \textbf{blocca} & $\geq 5{,}5$~dB & $0\,\%$ \\
\bottomrule
\end{tabular}
\end{table*}

\begin{figure*}[htbp]
\forcerectofloat
\centering
\includegraphics[width=\linewidth]{figures/fig15_attenuazione_materiali}
\caption{Attenuazione per materiale sui due radar, pannello a 20~cm dal sensore. A:
\ldb\ a 3~m, riduzione dell'energia in movimento rispetto alla baseline mediata su sei
trial. L'energia è un indice 0--100 e non una potenza, quindi vale la graduatoria e
non la conversione in decibel. B: \ldc\ a 1~m, energie grezze a 16~bit in decibel. Il
cartongesso (0,7~dB) attenua come plastica e cartone, legno e metallo sono limiti
inferiori.}
\label{fig:attenuazione}

\end{figure*}

Il \ldb\ rileva la presenza nel 100\,\% dei campioni attraverso tutti i materiali
dielettrici, con attenuazioni fra il 7 e il 42\,\% dell'energia. Il PIR è azzerato da
ognuno di essi. Il caso del vetro è il più immediato. È un materiale trasparente alla
luce visibile che il radar attraversa e che per il PIR è opaco, perché il vetro comune
non trasmette oltre i 4--5~\textmu m mentre l'emissione del corpo è centrata sui 10.
La graduatoria è identica sui due radar, a due distanze e in due unità (plastica,
cartone e vetroresina, vetro, legno, metallo). Il cartongesso, il materiale delle
pareti di un'aula, attenua come plastica e cartone. Con il metallo il \ldb\ riporta un
bersaglio a 30,0~cm con dispersione nulla, cioè il pannello stesso, e del soggetto non
resta traccia. Questo prova anche che l'aggiramento del pannello è sotto soglia. Le
percentuali del \ldb\ sono variazioni di un indice normalizzato e non si convertono in
decibel (\cref{sec:energia-distanza}). Sul \ldc\ la frazione di campioni sopra la
soglia di mantenimento del gate~2 passa dal 25--33\,\% senza ostacolo allo 0--1\,\%
dietro 10~mm di legno. Quell'esemplare non acquisirebbe quindi una persona a 1~m
dietro un pannello di legno, dove il \ldb\ a 3~m è al 100\,\%.

\paragraph{Portata residua con l'ostacolo.} Le percentuali non rispondono alla domanda
pratica del relatore. Se il radar arriva a 5~m senza ostacolo, dove arriva con una
porta di mezzo? Il \ldb\ ha quindi ripetuto la geometria percorrendo le distanze da 1
a 5~m con il cartongesso, e a 3--5~m con legno, vetro e una porta interna chiusa
(tamburata, senza vetro, con il sensore all'interno a 20--30~cm dal battente), tre
trial per punto (\cref{fig:portata-ostacoli}). Nessun dielettrico accorcia la portata
entro i 5~m della stanza. La presenza è al 100\,\% e la distanza è corretta in tutti i
39 trial con ostacolo, porta compresa. Ciò che degrada è il canale in movimento. La
frazione di campioni in cui è attivo a 5~m passa dal 78\,\% senza ostacolo al 44 con
il cartongesso, 26 con il vetro, 19 con il legno e 3\,\% con la porta. La presenza la
tiene il canale stazionario. La porta è trasparente a 3~m (energia pari alla baseline)
e quasi opaca a 5, più del legno pieno. Il fenomeno non è spiegato. Il PIR a 1~m passa
dal 44--97\,\% senza ostacolo allo 0\,\% con il cartongesso.

\begin{figure*}[htbp]
\forceversofloat
\centering
\includegraphics[width=\linewidth]{figures/fig14_portata_ostacoli}
\vspace{2.3cm}% spazio per la didascalia nel margine (tufte non lo riserva)
\caption{\ldb, portata residua con l'ostacolo a 20~cm dal sensore. A: energia in
movimento con cartongesso a 1--5~m e con legno, vetro e porta chiusa a 3--5~m, con
presenza al 100\,\% e distanza corretta in tutti i 39 trial. B: ciò che degrada è il
canale in movimento, la presenza la tiene il canale stazionario.}
\label{fig:portata-ostacoli}

\end{figure*}

Per il progetto il sensore può essere incassato nell'arredo, dietro legno o
cartongesso o dentro un guscio di vetroresina, senza perdere portata entro i 5~m.
Continua a rilevare anche attraverso una porta interna chiusa. Il PIR deve affacciarsi
direttamente sull'ambiente. Il metallo resta l'unico vincolo assoluto, e nel banco
reale la lamiera di rinforzo sta dietro al modulo e non davanti.

\section{Portata massima in corridoio}
\label{sec:portata-massima}

A 5~m nessuno dei tre sensori aveva mostrato il proprio limite in stanza. La prova di
portata è stata quindi svolta per ultima in un corridoio largo 120~cm
(\cref{fig:geometrie}d), con i sensori a 80~cm da terra e i punti a 6, 7 e 8~m visti
attraverso due vani porta. Venti minuti di corridoio vuoto non producono alcun evento
su radar e PIR. Il fondo dei gate lontani del \ldb\ ha massimo 7 contro una soglia di
15, quindi vani e muro di fondo non compaiono. I risultati dei tre sensori sono
riassunti in \cref{fig:portata-tre-sensori}.

\begin{figure*}[htbp]
\forcerectofloat
\centering
\includegraphics[width=\linewidth]{figures/fig17_portata_tre_sensori}
\caption{Portata dei tre sensori sulla stessa scala, stanza (agosto) e corridoio
(settembre). \ldb: 100\,\% fino a 6~m, zero a 7 per il tetto configurato. PIR a metà
corsa: attraversamento a circa 1~m/s rilevato fino a 5~m, quasi nullo a 6. Al massimo
di sensibilità 100\,\% a 5--6~m, 57\,\% a 7, 15\,\% a 8. \ldc\ (esemplare): si accende
entro 1~m (triangoli pieni) e mantiene la presenza fino a 2~m (triangoli vuoti).}
\label{fig:portata-tre-sensori}

\end{figure*}

\paragraph{\ldb.} Con un'oscillazione trasversale sul posto di circa 50~cm il radar
rileva al 100\,\% a 5 e 6~m e allo 0\,\% a 7~m in tre trial su tre. A 6~m la distanza
riportata ha massimo esattamente 600~cm in cinque trial su cinque, con dispersione di
2--5~cm contro i 10--11 di 5~m. Non è una misura più precisa, è il tetto degli otto
gate da 75~cm configurati. La presenza a 6~m è retta dal canale stazionario (energia
62--82). A 7~m i gate 7 e 8 restano molto sotto le soglie. Il tetto di 600~cm
dichiarato dal manuale è misurato. Il punto a 6~m non entra nella regressione della
distanza perché saturo.

\paragraph{PIR.} A metà corsa, la sensibilità di tutta la campagna, l'oscillazione
trasversale è rilevata al 24--33\,\% anche a 7~m, dove il radar per configurazione non
vede nulla. L'attraversamento vero a circa 1~m/s dà 0, 6 e 0\,\% a 6~m contro il
98,7\,\% di 5~m. Un transito più veloce alla stessa distanza dà il 30\,\% con cinque
impulsi. Al limite della portata la velocità del transito decide, al pari del tipo e
dell'ampiezza del movimento. Con il trimmer al massimo l'attraversamento è rilevato al
100, 100, $56{,}8 \pm 18{,}0$ e $14{,}7 \pm 10{,}0\,\%$ a 5, 6, 7 e 8~m. Il limite sta
fra 7 e 8~m, coerente con i 3--7~m del datasheet. Oltre il tetto del radar il PIR al
massimo vede ancora, ma solo chi attraversa il suo campo abbastanza in fretta. La
serie al massimo è dichiaratamente separata dal resto della campagna.

\paragraph{\ldc.} Con il gate massimo portato a 12 (840~cm), il valore di fabbrica
che copre gli 8~m dichiarati, l'esemplare acquisisce a 1~m e mantiene a 2~m. Da 3~m in
su le sue energie sono il fondo del corridoio vuoto gate per gate. Il gate massimo non
ha aggiunto nulla. Il limite è dell'esemplare e non della configurazione né della
stanza.

Le tre portate nella stessa geometria sono quindi queste. Il \ldb\ arriva a 6~m, per
tetto configurato e non per sensibilità. Il PIR arriva a 5--6~m a metà corsa e a 7--8~m
al massimo, sull'attraversamento. L'esemplare di \ldc\ acquisisce a 1~m e mantiene a
2~m. Per un sensore sotto un banco 6~m sono molto più del necessario. La domanda di
progetto va nella direzione opposta, come restringere il campo (\cref{sec:angolare}).

\section{Il secondo radar: che cosa cambia}
\label{sec:ld2420-risultati}
\label{sec:ld2420-taratura}
\label{sec:ld2420-fp}
\label{sec:ld2420-uprise}
\label{sec:ld2420-limiti}
\label{sec:ld2420-latenze}
\label{sec:ld2420-angolare}

La campagna sul secondo radar è stata svolta per intero entro i 2~m di portata
dell'esemplare, con le deroghe alla parità di configurazione dichiarate nella
\cref{sec:protocollo-ld2420}, per 161 trial validi. I suoi risultati test per test
sono già nelle figure delle sezioni precedenti. Questa sezione raccoglie ciò che
cambia rispetto al \ldb. Ogni numero è attribuito all'esemplare, mai al modello, e la
\cref{tab:ld2420-sintesi} li riassume.

\paragraph{La taratura è una condizione necessaria.} Con le soglie di fabbrica e il
gate massimo 12 il modulo dichiara presenza permanente nella stanza di prova, anche
vuota. Il muro a circa 5~m sta nel gate~7 e supera la soglia di mantenimento nel
4\,\% dei campioni, e con un ritardo di scomparsa di 30~s basta un superamento ogni
mezzo minuto. Con il gate massimo a 6 il muro esce dal campo. La scansione del rumore
dello strumento del produttore non basta, perché sottostima le code in modo
ripetibile. La soglia di mantenimento del gate~5 esce a 13,8~dB in tre scansioni,
contro un fondo misurato su 120~s con media 15 e massimo 34. Le soglie definitive
applicano la formula del manuale al fondo misurato. Il \ldb\ nella stessa stanza
funziona a soglie di fabbrica con zero falsi positivi in 6,55~h. Per centinaia di
banchi è una differenza operativa sostanziale.

\paragraph{Falsi positivi.} Le 181 riaccensioni in 7~h di \cref{sec:falsi-positivi}
hanno una firma precisa, episodi di 31,8~s mediani, cioè un trigger isolato più i 30~s
di ritardo. L'attribuzione per gate indica il gate~0 in 80 casi su 181, e il gate~0 è
l'accoppiamento fra trasmettitore e ricevitore, non una distanza. Portando il solo
gate~0 sopra la coda del rumore, un'ora di giorno dà 5 riaccensioni, 5,4 eventi/h. Il
residuo viene dai gate 1 e 2, che non si possono alzare perché sono quelli in cui la
persona a 50~cm e a 1~m produce il suo segnale con margine quasi nullo. Su questo
esemplare rumore e persona si sovrappongono, e 5,4 eventi/h è un pavimento
strutturale. Il modulo inoltre stima il fondo all'accensione. Un corpo a 20~cm nei
primi 60~s lo lascia in presenza permanente senza risposta alla persona, finché non
viene spento per un quarto d'ora. Il \ldb\ non l'ha mai fatto in decine di riavvii.

\paragraph{Nessuna grandezza continua.} La dose-risposta di \cref{sec:dose-risposta},
ripetuta sul \ldc\ nella stessa geometria, non gradua il movimento. L'energia del
gate~2 vale 13,4 a stanza vuota, $41{,}0 \pm 9{,}3$ da immobile, $34{,}8 \pm 1{,}0$
con i micro-movimenti e 31,0 in cammino. L'ordine è rovesciato rispetto al \ldb\
(68,6, 84,4, 99,3) e le condizioni si sovrappongono. Il canale separa vuoto da
occupato e poi non distingue. Su questo esemplare non esiste quindi la grandezza da
cui costruire l'indice di vitalità.

\paragraph{Acquisizione e gate minimo.} In venti trial di cammino sul posto da 0,5
a 2~m la presenza è al 100\,\%, ma il modulo acquisisce fino a circa 1~m e da lì
mantiene soltanto. La distanza riportata è discussa in \cref{sec:distanza} e le due
persone in \cref{sec:due-persone}. Il gate minimo è l'unica funzione che il \ldc\ ha e
il \ldb\ no. Non tocca la decisione di presenza (100\,\% con la persona nel gate
escluso) e confina la distanza al bordo della finestra.

\paragraph{Latenze e rilascio.} Le misure sono in \cref{sec:latenza}. Il rilascio
lento osservato all'inizio della campagna (87--120~s) era il parametro di fabbrica di
30~s, non un'isteresi dell'esemplare. La prova con il ritardo a 5~s ha chiuso anche
l'unità del parametro, su cui manuale e documento di protocollo si contraddicono.
Sono secondi.

\paragraph{Respiro.} Con il metronomo a 15 atti al minuto a 1~m, il criterio del
\cref{sec:respiro} recupera la frequenza in un trial su tre (due su tre a gate
fissato), con controlli negativi puliti, contro 13 su 15 del \ldb. La persona occupa
due gate soli contro i 7--14 canali concordi del \ldb, e il rapporto segnale-rumore
massimo è 6,1 contro 12,8. La saturazione è zero, ma il vantaggio dei 16~bit è
annullato dal segnale debole dell'esemplare.

\begin{table*}[htbp]
\forcerectofloat
\centering
\caption{I due radar a confronto, tutte grandezze misurate in questa tesi. Le righe
del \ldc\ si riferiscono all'esemplare in dotazione, con soglie tarate.}
\label{tab:ld2420-sintesi}

\small
\begin{tabularx}{\linewidth}{@{}Xll@{}}
\toprule
\textbf{Grandezza} & \textbf{\ldb} & \textbf{\ldc\ (esemplare)} \\
\midrule
Funziona a soglie di fabbrica nella stanza di prova & sì & no, presenza permanente \\
Falsi positivi a stanza vuota & $\leq 0{,}43$/h (0 in 6,6~h) & 26,2/h tarato, 5,4/h con gate~0 alzato \\
Portata su persona in movimento & 6~m (tetto configurato) & 1~m acquisizione, 2~m mantenimento \\
Persona immobile a 1~m / sotto il banco & $100\,\%$ / $100\,\%$ & $100\,\%$ / $98{,}4\,\%$ \\
Distanza su bersaglio fermo & sì, dispersione 1,5~cm & no (manuale \S8, misurato) \\
Distanza su bersaglio in moto & continua fino a 5~m, $R^2 = 0{,}99965$ & valida fino a 1~m, poi stantia \\
Graduazione del movimento & 68,6, 84,4, 99,3 & piatta e rovesciata \\
Latenza d'ingresso & $5{,}36 \pm 0{,}30$~s & $7{,}28 \pm 1{,}40$~s \\
Rilascio (parametro 5~s) & $18{,}4$~s & $7{,}7$~s (campo più corto) \\
Copertura angolare a 1~m & $\geq \pm 90$\textdegree & $\approx \pm 75$\textdegree \\
Cartongesso 10~mm & $-26\,\%$ a 3~m, portata intatta & 0,7~dB \\
Legno 10~mm & $-42\,\%$, portata intatta & non acquisirebbe a 1~m \\
Respiro a ritmo imposto & 13/15 trial & 1/3 (2/3 a gate fisso) \\
Due persone & un bersaglio per canale & un bit \\
Comportamento all'avvio & stabile in decine di riavvii & fondo corrotto da un corpo vicino \\
\bottomrule
\end{tabularx}
\end{table*}

\paragraph{Un riscontro indipendente.}
\label{sec:ld2420-riscontro}
Una relazione di progetto svolta a Camerino nello stesso anno accademico arriva a
conclusioni sovrapponibili, con un altro esemplare di ciascun modulo, un'altra libreria
e un altro protocollo sperimentale~\cite{dellaquila2026}. Sul \ldc\ riporta un flag di
presenza inaffidabile, un primo gate con energie alte anche a stanza vuota ed escluso
dalla logica, la necessità di tarare le soglie per gate e, anche dopo, la
sovrapposizione fra stanza vuota e soggetto poco attivo. Sul \ldb\ non osserva falsi
positivi con finestre aperte, tende mosse dal vento e in ambiente esterno, e giudica
stabile la presenza per tutta la permanenza del soggetto. Le differenze sono due. La
relazione giudica buona l'accuratezza di distanza del \ldc, coerente con
l'attribuzione all'esemplare del nostro limite di portata. Giudica inoltre il \ldb\
inadatto al respiro, ma legge la distanza con il modulo sul torace, mentre qui si
legge l'energia per gate a 2~m.

\paragraph{La lettura per il progetto.} Nello scenario DIPME il \ldc\ rileva la
persona immobile sotto il banco, e questo va detto. Ma dice se e non dove, non gradua
il movimento, richiede una taratura per installazione e conserva 5,4 riaccensioni
l'ora. Il \ldb\ funziona a soglie di fabbrica con zero falsi positivi, riporta la
distanza anche sul bersaglio fermo e fornisce la grandezza continua da cui nasce
l'indice di vitalità. Un secondo esemplare di \ldc\ permetterebbe di separare i limiti
del modello da quelli del pezzo.

\section{Sintesi del confronto}
\label{sec:discussione}
\label{sec:limiti}

La \cref{tab:sintesi} raccoglie in una sola vista i risultati del capitolo. Tutte le
righe con soggetto in movimento o con micro-movimenti sono acquisite con il PIR nella
configurazione a lui più favorevole.

\begin{table*}[htbp]
\forceversofloat
\centering
\caption{Sintesi della campagna sperimentale. ``Rilev.'' è la frazione di campioni in
cui il sensore riporta la presenza, con la ground truth indicata. Le ultime due righe
si riferiscono al \ldc\ (esemplare, soglie tarate), tutte le altre al \ldb.}
\label{tab:sintesi}

\small
\begin{tabularx}{\linewidth}{@{}llrrX@{}}
\toprule
\textbf{Scenario} & \textbf{Verità} & \textbf{Rilev.\ PIR} & \textbf{Rilev.\ radar} &
\textbf{Riferimento} \\
\midrule
Stanza vuota, 7,05~h & assente & $0$ ev/h & $0$ ev/h & \cref{sec:falsi-positivi} \\
Fermo seduto, 2,3~m & presente & $0{,}22\,\%$ & $\mathbf{100\,\%}$ & \cref{sec:persona-ferma} \\
Fermo in piedi, 1~m & presente & $1{,}52\,\%$ & $\mathbf{100\,\%}$ & \cref{sec:dose-risposta} \\
Micro-movimenti, 1~m & presente & $51{,}60\,\%$ & $\mathbf{100\,\%}$ & \cref{sec:dose-risposta} \\
Cammino sul posto, 1~m & presente & $85{,}20\,\%$ & $\mathbf{100\,\%}$ & \cref{sec:dose-risposta} \\
Cammino sul posto, 2--5~m & presente & $0{,}00\,\%$ & $\mathbf{100\,\%}$ & \cref{sec:dose-risposta} \\
Attraversamento 2--5~m & presente & $99$--$100\,\%$ & $\mathbf{100\,\%}$ & \cref{sec:dose-risposta} \\
Sotto il banco, immobile & presente & $1{,}32\,\%$ & $\mathbf{100\,\%}$ & \cref{sec:sotto-banco} \\
Sotto il banco, micro-mov. & presente & $96{,}52\,\%$ & $\mathbf{100\,\%}$ & \cref{sec:sotto-banco} \\
Due persone, ogni geometria & presenti & --- & $\mathbf{100\,\%}$ & \cref{sec:due-persone} \\
Vicino a 3~m, gate 2 & assente & $0{,}00\,\%$ & $0{,}00\,\%$ & \cref{sec:angolare} \\
Vicino fermo a 90\textdegree, gate 2 & assente & $0{,}00\,\%$ & $0{,}00\,\%$ & \cref{sec:angolare} \\
Vicino a 90\textdegree\ in movimento & assente & $12\,\%$ & $\mathbf{100\,\%}$ & \cref{sec:angolare} \\
Dietro cartongesso, legno, vetro, porta, 3--5~m & presente & $0\,\%$ (1~m) & $\mathbf{100\,\%}$ & \cref{sec:ostacoli} \\
Corridoio 6~m / 7~m, oscillazione & presente & $12$--$51$ / $24$--$33\,\%$ & $\mathbf{100\,\%}$ / $0\,\%$ & \cref{sec:portata-massima} \\
Attraversamento 6~m, PIR al massimo & presente & $\mathbf{100\,\%}$ & $100\,\%$ & \cref{sec:portata-massima} \\
Stanza vuota, \ldc\ tarato, 7~h & assente & --- & 26,2 ev/h & \cref{sec:ld2420-risultati} \\
Sotto il banco, immobile, \ldc & presente & --- & $98{,}4\,\%$ & \cref{sec:sotto-banco} \\
\bottomrule
\end{tabularx}
\end{table*}

\Needspace{6\baselineskip}
\paragraph{Dove il mmWave vince.}
\begin{itemize}[leftmargin=1.4em,itemsep=2pt]
  \item Sulla persona ferma il divario è categorico, non graduale. A 2,3~m seduto,
        a 1~m in piedi e sotto il banco a 60~cm il PIR ha il 98,5--99,8\,\% di falsi
        negativi e il radar lo 0,00\,\%. Tre geometrie, lo stesso esito.
  \item Il divario è attribuibile a una sola variabile, l'immobilità del soggetto,
        perché la dose-risposta la isola a geometria, distanza e configurazione
        costanti.
  \item Il radar è più veloce anche dove il PIR dovrebbe essere avvantaggiato,
        all'ingresso da una porta, in dieci trial su dieci.
  \item Il radar aggiunge informazione, non solo affidabilità. La distanza è lineare
        con $R^2 = 0{,}99965$. L'energia e la dispersione della distanza sono
        indicatori continui della quantità di movimento.
  \item Sui falsi positivi i due sensori sono equivalenti ($\leq 0{,}43$ eventi/h).
        Il radar attraversa cartongesso, legno, vetro e una porta chiusa senza perdere
        portata, mentre il PIR è azzerato da ogni materiale.
\end{itemize}

\paragraph{Dove il mmWave non vince.}
\begin{itemize}[leftmargin=1.4em,itemsep=2pt]
  \item Il consumo, tre ordini di grandezza in più (80~mA contro 50~\textmu A), è il
        vincolo che determina l'architettura del \cref{cap:conclusioni} e che il
        \cref{cap:consumi} quantifica.
  \item Il rilascio è più lento del PIR di 5,4~s e più lento del parametro
        dichiarato (coda di circa 9~s contro 5~s configurati).
  \item A distanza ravvicinata i canali più informativi saturano e i gate stazionari
        non sono disponibili sotto 1,5~m. Sotto l'arredo la localizzazione è dieci
        volte più rumorosa. La presenza resta al 100\,\%.
  \item Riporta un bersaglio per canale, quindi non è un contatore di presenze.
        Richiede UART, protocollo a frame ed elaborazione numerica contro il piedino
        digitale del PIR.
  \item Il modello non è il modulo. Il secondo radar provato, con portata dichiarata
        maggiore e consumo minore, sull'esemplare in dotazione vede fino a 2~m, non
        rilascia a soglie di fabbrica e conserva 5,4 falsi positivi l'ora dopo la
        taratura. I vantaggi misurati sono del \ldb, non della categoria mmWave in
        astratto.
\end{itemize}

\paragraph{Limiti del setup.} Tutti i risultati comparativi sono di un solo soggetto,
in una stanza domestica e non in un'aula, a 25--29\,\textdegree C\ (sfavorevoli al PIR
in rilevamento). C'è un solo esemplare per modello e nessuna strumentazione di misura
indipendente per distanze e consumi. Gli scenari portanti hanno cinque o dieci trial,
quelli di selettività e a due persone soltanto tre. Su tre trial le affermazioni sono
formulate come assenza di effetto osservato. La prima parte della campagna è in
posizione L del PIR. Le serie sensibili sono state ripetute in H, e i due gruppi non
sono uniti nella regressione della distanza perché separati da un rimontaggio del
setup. Nessuno di questi limiti tocca il risultato centrale, che è robusto per la
natura del divario misurato. La cecità del PIR sulla persona ferma è una conseguenza
del principio fisico. I dati la confermano in tre geometrie indipendenti, con cinque
ripetizioni ciascuna, nella configurazione più favorevole al PIR.
